考虑一台无磁场、满足微观时间反演对称性的两端热电发电机。两热库温度为 \(T_h=T+\Delta T/2\)、\(T_c=T-\Delta T/2\)，其中 \(T>0\)、\(0<\Delta T\ll T\)。热流 \(J_q>0\) 从热端流向冷端；负载电压 \(V>0\) 和电流 \(I>0\) 对应向负载输出功率 \(P=IV\)。仅保留电流对热力学力的一阶响应、熵产生的二阶项。定义
\[
X_1=-\frac{V}{T},\qquad X_2=\frac{\Delta T}{T^2}>0,\qquad
\binom{I}{J_q}=
\begin{pmatrix}L_{11}&L_{12}\\L_{12}&L_{22}\end{pmatrix}
\binom{X_1}{X_2}.
\]
假定 \(L_{11},L_{22},L_{12}>0\) 且 \(\det L>0\)；以 \(\eta_{\mathrm C}^{(1)}=\Delta T/T=T X_2\) 表示线性响应阶的卡诺效率。这里 \(J_q\) 是按参考温度 \(T\) 定义的热流，因而输出效率 \(P/J_q\) 只要求在线性响应首个非零阶精确。
(1) 从热库熵变及电功推出熵产生率 \(\dot S_{\rm i}=I X_1+J_q X_2\)，严格给出二次型非负条件。引入耦合强度 \(q=L_{12}/\sqrt{L_{11}L_{22}}\)，证明 \(0<q<1\)，并完成平方以说明不可逆热漏与电流松弛两种耗散来源。
(2) 将负载参数化为 \(X_1=-x(L_{12}/L_{11})X_2\)，其中 \(0\le x\le1\)。求 \(I,J_q,P,\dot S_{\rm i}\) 和归一化效率 \(\eta/\eta_{\mathrm C}^{(1)}\)，确定短路、开路边界，证明发电区热流为正且输出功不超过线性响应可用热功。
(3) 分别在 \(0\le x\le1\) 上求最大输出功率及对应效率、最大效率及其唯一最优负载。引入 \(ZT=q^2/(1-q^2)\) 化简最大效率；讨论 \(q\to0\) 与 \(q\to1^-\) 时的功率、效率和可逆性，不能将零功率极限误写成有限功率卡诺机。
(4) 对一根等温参数 \(G>0\)（电导）、\(S>0\)（Seebeck系数）、\(K>0\)（零电流热导）的均匀热电腿，给定本构关系
\[
I=G(S\Delta T-V),\qquad J_q=STI+K\Delta T .
\]
从该关系独立重建 \(L_{ij}\)，证明 \(ZT=GS^2T/K\) 及 Kelvin 关系 \(\Pi=ST\) 的作用。求负载电阻 \(R_L=V/I\) 的最大功率匹配与最大效率匹配；直接从物理本构关系验证熵产生、开路热漏及 \(K\to0^+\) 极限，并核对量纲。